%-----------------------------------LICENSE------------------------------------%
%   This file is part of Mathematics-and-Physics.                              %
%                                                                              %
%   Mathematics-and-Physics is free software: you can redistribute it and/or   %
%   modify it under the terms of the GNU General Public License as             %
%   published by the Free Software Foundation, either version 3 of the         %
%   License, or (at your option) any later version.                            %
%                                                                              %
%   Mathematics-and-Physics is distributed in the hope that it will be useful, %
%   but WITHOUT ANY WARRANTY; without even the implied warranty of             %
%   MERCHANTABILITY or FITNESS FOR A PARTICULAR PURPOSE.  See the              %
%   GNU General Public License for more details.                               %
%                                                                              %
%   You should have received a copy of the GNU General Public License along    %
%   with Mathematics-and-Physics.  If not, see <https://www.gnu.org/licenses/>.%
%------------------------------------------------------------------------------%
\documentclass{article}
\usepackage{amssymb}
\usepackage{hyperref}
\hypersetup{colorlinks = true}
\title{Completeness and Complex Numbers}
\author{MIT 18.100P}
\date{\today}
\setlength{\parindent}{0em}
\setlength{\parskip}{0em}
\begin{document}
    \maketitle
    Completeness is one of the most fundamental properties of the real
    numbers. It is the lack of completeness in $\mathbb{Q}$ that led us to
    construct $\mathbb{R}$ in the first place. In class we discussed that once
    we have $\mathbb{R}$, we can construct $\mathbb{C}$ using the Cartesian
    product, and there is only one \textit{nice} arithmetic on $\mathbb{C}$.
    That is, there is only set of operations $+$ and $\cdot$
    (up to \textit{isomorphism}, whatever that fancy word means) such that:
    \begin{enumerate}
        \item
            $+$ is just vector addition in $\mathbb{R}^{2}$.
        \item
            $\cdot$ distributes over $+$.
        \item
            $\cdot$ is compatibile with the real numbers, meaning
            \begin{equation}
                (a,\,0)\cdot(c,\,d)
                =
                (ac,\,ad).
            \end{equation}
            We treat $(a,\,0)$, which lies
            along the $x$ axis, the same way we treat the real number $a$.
        \item
            Division by non-zero complex numbers is possible. That is,
            if $(a,\,b)\ne(0,\,0)$, then $(a,\,b)^{-1}$ is well defined.
        \item
            $\cdot$ is compatibile with Euclidean geometry, meaning:
            \begin{equation}
                \|(a,\,b)\cdot(c,\,d)\|
                =
                \|(a,\,b)\|
                \;
                \|(c,\,d)\|
            \end{equation}
    \end{enumerate}
    To do this, $\cdot$ must be:
    \begin{equation}
        (a,\,b)\cdot(c,\,d)
        =
        (ac-bd,\,ad+bc)
    \end{equation}
    Perhaps more familiar, we write $1=(1,\,0)$ and $i=(0,\,1)$, and then have:
    \begin{equation}
        (a,\,b)\cdot(c,\,d)
        =
        (a+ib)\cdot(c+id)
        =
        (ac-bd)+i(ad+bc)
    \end{equation}
    With this, $(\mathbb{C},\,+,\,\cdot)$ forms a \textit{field},
    just like $(\mathbb{Q},\,+,\,\cdot)$ and $(\mathbb{R},\,+,\,\cdot)$.
    There is also a metric on $\mathbb{C}$ given by the Pythagorean formula:
    \begin{equation}
        d\Big((a,\,b),\,(c,\,d)\Big)
        =
        \sqrt{(c-a)^{2}+(d-b)^{2}}
    \end{equation}
    Your goal is to prove that $\mathbb{C}$ is complete, just like $\mathbb{R}$.
    You should discuss:
    \begin{enumerate}
        \item
            Sequences.
        \item
            Convergence and limits.
        \item
            Cauchy sequences.
    \end{enumerate}
    You should also mention as a necessary lemma (but do not prove)
    that $\mathbb{R}$ is complete.
    \par\hfill\par
    \textbf{\Large{Guidelines}}
    \par\hfill\par
    Your article must be written in \LaTeX{} or \TeX, the primary typesetting
    languages for mathematics and other sciences. \LaTeX{} is preferred since it
    is much easier to hit the ground running and start writing documents.
    \begin{itemize}
        \item
            Your article is expected to be at least 3 pages,
            but realistically something in the realm of 5-7 is more likely.
        \item
            50\% of the grade is for \textit{mathematical correctness}.
            This addresses the following:
            \begin{itemize}
                \item
                    Are you using definitions correctly?
                \item
                    Are you applying axioms and other theorems
                    appropriately in your arguments?
                \item
                    Is your proof valid?
                \item
                    If you are using figures, are they relevant and explained
                    properly?
            \end{itemize}
        \item
            50\% comes from \textit{writing style}:
            \begin{itemize}
                \item
                    Does the paper meet all the expectations
                    from writing assignment 1, i.e.:
                    \begin{itemize}
                        \item
                            Does the paper have an introduction that identifies
                            and signals the purpose and the main result for a
                            general audience (i.e. not someone currently
                            enrolled in 18.100P or necessarily currently
                            taking real analysis)?
                        \item
                            Does the paper set up key definitions and give the
                            intuition of the approach used to prove the result
                            so that the reader can point to the places in the
                            text where these are given?
                        \item
                            Does the paper use of statement environments
                            and proof environments?
                        \item
                            Does the paper write in paragraphs,
                            integrating words and symbols?
                    \end{itemize}
                \item
                    Does the paper guide the reader through the ideas of the
                    paper, signposting the structure, motivating sections, and
                    motivating definitions and lemmas?
                \item
                    Does the paper guide the reader through the reasoning of
                    the proofs by giving overviews of the proof strategies,
                    justification for intermediate claims,
                    and identifying key steps in calculations?
                \item
                    Does the paper include at least one figure referred to and
                    discussed in the text (and no figures that don't have the
                    above)?
            \end{itemize}
    \end{itemize}
    \par\hfill\par
    \textbf{\Large{Deadlines}}
    \par\hfill\par
    \textbf{Rough Draft: 2025/03/13}
    \par
    A full draft of your article.
    You will receive feedback on your writing shortly after break.
    \par\hfill\par
    \textbf{Peer Review: 2025/03/27}
    \par
    You will participate in a peer review session with your classmates.
    \par\hfill\par
    \textbf{Final Due: 2025/04/04}
    \par
    After making revisions, you will turn in your final draft for grading.
\end{document}
